\title{Controlling Text-to-Image Diffusion by Orthogonal Finetuning}

\begin{abstract}
Text-to-image diffusion models have demonstrated remarkable proficiency in producing photorealistic images conditioned on textual prompts. A central and unresolved issue, however, is how to reliably steer these models toward specific downstream objectives. In response to this, we propose a systematic approach called Orthogonal Finetuning~(OFT) for modifying text-to-image diffusion models for a variety of downstream applications. In contrast to prior work, OFT mathematically guarantees the conservation of hyperspherical energy, a metric describing the pairwise interactions between neurons on the unit hypersphere. We establish that maintaining this property is essential to uphold the models' semantic image synthesis capacity. To further enhance the robustness of the finetuning process, we introduce Constrained Orthogonal Finetuning (COFT), which enforces an explicit radial constraint within the hyperspherical representation. Our study focuses on two key scenarios: subject-driven generation, where the system must render instances of a target subject based on a limited number of input images and accompanying text prompt; and controllable generation, where the model is expected to synthesize images that obey supplementary control cues. Empirical evidence supports that our OFT framework not only achieves superior image quality but also accelerates convergence when compared to contemporary alternatives.
\end{abstract}

\section{Introduction}

Modern text-to-image diffusion architectures~\cite{saharia2022photorealistic,ramesh2022hierarchical,rombach2022high} set new standards for synthesizing high-fidelity visuals in response to textual descriptions. Despite their strengths, the guidance offered solely by text often lacks the specificity and granularity necessary for fine-level manipulation of the generated content. In this work, we aim to overcome these limitations by concentrating on two particular categories of text-to-image tasks:

\begin{itemize}[leftmargin=*,nosep]

    \item \textbf{Subject-driven generation}~\cite{ruiz2023dreambooth}: Using only a small collection of reference images for a given subject, the model is tasked with creating new images of that subject placed in varied scenarios, under the direction of a text description.
    \item \textbf{Controllable generation}~\cite{zhang2023adding,mou2023t2i}: Provided with additional input constraints such as canny edges or segmentation masks, the goal is to produce images that conform to both the control cue and the text-based prompt.
\end{itemize}

At their core, both tasks concern the challenge of finetuning text-to-image diffusion models in a manner that maintains the generative capabilities acquired during pretraining. We delineate the objectives of a robust finetuning approach as follows: (1) \emph{training efficiency}, characterized by a reduction in both the number of adjustable parameters and the epochs required; and (2) \emph{generalizability preservation}, referring to the retention of the model’s capacity for high-fidelity and diverse generation. Traditionally, finetuning approaches either adjust existing neuron weights with minimal learning rates (\emph{e.g.}, \cite{ruiz2023dreambooth}) or augment the model with small, re-parameterized modules (\emph{e.g.}, \cite{hulora2022,zhang2023adding}). However, these straightforward techniques are insufficient to reliably ensure that generative performance from pretraining is conserved. Furthermore, there remains limited theoretical guidance for developing optimal finetuning methodologies and for determining critical hyperparameters, such as training duration. Central to these issues is the absence of an explicit metric for evaluating how well a model’s pretrained generative abilities are retained. Current strategies typically rest on the tacit premise that a reduced Euclidean distance between the weights of the finetuned and pretrained models correlates with superior preservation of generative function, thus often employing minimal learning rates. While this might occasionally yield positive results, we contend that Euclidean proximity alone does not comprehensively reflect the preservation of semantic integrity. Consequently, adopting a more structurally informed metric to compare the pretrained and finetuned models would substantially improve both the retention of generative performance and the stability of the finetuning process.

Building on empirical findings that hyperspherical similarity effectively captures semantic relationships~\cite{liu2017deep,liu2018decoupled,chen2020angular}, our approach leverages hyperspherical energy~\cite{liu2018learning} to represent the pairwise structural relationships between neurons. Specifically, hyperspherical energy, which aggregates the hyperspherical similarities (such as cosine similarity) across all neuron pairs within a single layer, reflects how evenly the neurons are distributed on the unit hypersphere~\cite{liu2021learning}. We propose that an optimal finetuned model should closely match the hyperspherical energy of its corresponding pretrained version. Although one could simply regularize the model to maintain constant hyperspherical energy throughout finetuning, this alone does not ensure effective minimization of energy discrepancies. To address this, we utilize a key property of hyperspherical energy: when an identical orthogonal transformation is applied across all neurons, the pairwise hyperspherical similarities remain unchanged. Drawing on this invariance, we introduce Orthogonal Finetuning~(OFT), a method for adapting large text-to-image diffusion models to target tasks while preserving hyperspherical energy. The principal mechanism involves learning an orthogonal transformation, shared across layers, that maintains inter-neuron angles. From another perspective, OFT can be interpreted as recalibrating the coordinate system for neurons within a layer. Additionally, acknowledging that closer Euclidean proximity between the finetuned and pretrained models tends to correlate with better retention of the pretrained capabilities, we develop an extension called Constrained Orthogonal Finetuning~(COFT). COFT restricts the finetuned neuron representations to remain within a hypersphere of fixed radius centered at the pretrained configuration.

The rationale behind the effectiveness of orthogonal transformations for neuron finetuning is partly rooted in insights from the 2D Fourier transform, wherein an image is represented through its magnitude and phase spectra. Notably, the phase spectrum—encoding the angular relationships between input and basis vectors—retains the essential semantic information. For instance, when reconstructing an image using its original phase spectrum in combination with an arbitrary magnitude spectrum, the semantics remain largely preserved. This observation indicates that adjusting neuron orientations is central to imparting semantic changes in the resulting image, which aligns with objectives in both subject-driven and controllable image generation tasks. However, unrestricted modification of neuron directions would likely compromise the generative capabilities inherited from pretraining. To manage this, we advocate for maintaining the angular relationships among all pairs of neurons, supporting the hypothesis that such inter-neuronal angles are fundamental to knowledge representation within neural networks. This intuition leads to the proposal of a layer-wise, shared orthogonal transformation across neurons in each layer, ensuring that the hyperspherical energy remains constant.

Our approach also builds on ideas from orthogonal over-parameterized training~\cite{liu2021orthogonal}, which involves training classification networks by applying orthogonal transformations to the network initialized at random. Random initializations, as demonstrated, typically have low expected hyperspherical energy, and the objective in \cite{liu2021orthogonal} is to sustain this low energy throughout training—a property known to facilitate improved generalization in classification tasks~\cite{liu2018learning,lin2020regularizing}. Their work demonstrates that orthogonal transformations are sufficiently expressive for learning generalized classification models. Distinct from their focus, our work concentrates on the finetuning of text-to-image diffusion models to enhance both controllability and generative quality in downstream applications. The distinction between OFT and the method in \cite{liu2021orthogonal} is twofold. Firstly, while \cite{liu2021orthogonal} targets the reduction of hyperspherical energy, OFT is designed to preserve the energy value inherent to the pretrained model, safeguarding the model’s latent semantic structures during finetuning—since indiscriminate minimization of hyperspherical energy may compromise these semantics. Secondly, OFT explicitly seeks to minimize the shift from the pretrained state, leading to a constrained approach, whereas no such regularization is present in \cite{liu2021orthogonal}. The challenge of effective finetuning thus lies in balancing adaptability with preservation, a trade-off that the OFT framework addresses with efficacy. The following points summarize our primary contributions:

\begin{itemize}[leftmargin=*,nosep]

    \item We introduce Orthogonal Finetuning, a new finetuning strategy that enhances the controllability of text-to-image diffusion models. To address stability concerns, we additionally present a constrained variant that restricts the angular shift from the pretrained model's parameters.
    \item In contrast to conventional finetuning methods, OFT enables model adaptation while rigorously maintaining hyperspherical energy—a property that, through our empirical investigations, proves crucial for retaining the generative semantic integrity of the original model.
    \item We employ OFT in both subject-driven and controllable generation contexts. Through a thorough empirical analysis, we establish that OFT achieves substantial gains over previous methods with respect to generation fidelity, convergence rate, and the robustness of the finetuning process. Notably, OFT is more sample-efficient, attaining strong convergence with significantly fewer training examples and epochs.
    \item To assess controllable generation, we propose a novel control consistency metric that quantifies controllability by estimating the control signal inferred from a generated image and comparing it against the initial control signal. This metric offers further evidence of OFT’s superior controllability.
\end{itemize}

\section{Related Work}

\textbf{Text-to-image diffusion models}. The field of text-to-image generation has witnessed remarkable advances~\cite{saharia2022photorealistic,ramesh2022hierarchical,rombach2022high,gu2022vector,nichol2022glide}, largely propelled by progress in diffusion-based generative approaches~\cite{ho2020denoising,song2021score,song2021denoising,dhariwal2021diffusion} and joint vision-language representation frameworks~\cite{su2020vlbert,lu2019vilbert,radford2021learning,wang2021simvlm,li2022blip,li2023blip,alayrac2022flamingo,schuhmann2022laion}. For instance, GLIDE~\cite{nichol2022glide} and Imagen~\cite{saharia2022photorealistic} operate in pixel space, with GLIDE jointly training its text encoder and diffusion prior using paired datasets, whereas Imagen incorporates a fixed, pretrained text encoder. Diffusion models working in latent spaces, like Stable Diffusion~\cite{rombach2022high} and DALL-E2~\cite{ramesh2022hierarchical}, have also achieved prominence: Stable Diffusion utilizes VQ-GAN~\cite{esser2021taming} for latent visual representation, while DALL-E2 leverages CLIP~\cite{radford2021learning} to form a shared image-text embedding space. Beyond diffusion paradigms, generative adversarial networks~\cite{reed2016generative,zhang2017stackgan,xu2018attngan,li2019controllable} and autoregressive strategies~\cite{ramesh2021zero,ding2021cogview,wu2022nuwa,yuscaling} have also been explored for text-to-image synthesis. OFT is designed as a model-agnostic finetuning technique, applicable across diverse text-to-image diffusion models.

\textbf{Subject-driven generation}. Approaches to mitigate subject alteration have included methods like \cite{avrahami2022blended,nichol2022glide}, which employ user-supplied masks as auxiliary constraints. Inversion-based techniques~\cite{choi2021ilvr,dhariwal2021diffusion,ramesh2022hierarchical,gal2022image} facilitate contextual edits while maintaining the subject unaltered, and \cite{hertz2022prompt} allows both localized and comprehensive manipulations without explicit masks. Nevertheless, these strategies often struggle to preserve fine-grained, identity-specific aspects of the subject. Pivotal Tuning~\cite{roich2022pivotal} addresses this by refining the generator near an inverted latent code, incorporating regularization to retain identity features. Similarly, \cite{nitzan2022mystyle} establishes personalized facial priors from an individual's image set, and \cite{casanova2021instance} enables the generation of varied renditions of an instance at the expense of losing certain instance-unique attributes. DreamBooth~\cite{ruiz2023dreambooth} customizes diffusion models via tailored tokens and a small set of subject images, using reconstruction and class-based prior preservation losses for finetuning. Building on the DreamBooth structure, OFT replaces conventional small-step finetuning with parameter updates constrained to orthogonal transformations.

\textbf{Controllable generation}. Image-to-image translation is often conceptualized as a controllable generation challenge, for which earlier work has predominantly employed conditional generative adversarial networks~\cite{isola2017image,zhu2017toward,wang2018high,park2019semantic,choi2018stargan}. Diffusion models have emerged as another approach for this task~\cite{saharia2022palette,voynov2022sketch,wang2022pretraining}. Recently, ControlNet~\cite{zhang2023adding} introduced a mechanism to enhance the controllability of pretrained diffusion models through targeted finetuning and incorporation of additional control inputs, achieving strong results. Similarly, T2I-Adapter~\cite{mou2023t2i} fine-tunes diffusion models to provide greater control over the generation process. Building on the frameworks introduced in \cite{zhang2023adding,mou2023t2i}, we employ OFT for the finetuning of pretrained diffusion models, demonstrating more robust controllability even when limited by smaller training sets and reduced finetuning parameters. Importantly, OFT incurs no extra computational cost during inference.

\textbf{Model finetuning}. Adapting large pretrained models to new downstream applications through finetuning has become standard practice~\cite{devlin2018bert,brown2020language,he2022masked}. Adaptation techniques (such as those in \cite{hulora2022,houlsby2019parameter,pfeiffer2020adapterfusion}) have been extensively researched in the field of natural language processing. LoRA~\cite{hulora2022}, closely related to OFT, operates under the assumption that weight updates during finetuning are of low rank. Distinct from this, OFT updates neuron weights through layer-shared orthogonal transformations, which are multiplicative rather than additive, and guarantees the preservation of pair-wise angles within a layer. This theoretical property leads to improved training stability.

\section{Orthogonal Finetuning}

\subsection{Why Does Orthogonal Transformation Make Sense?}

We begin our analysis by examining the rationale for employing orthogonal transformations in the finetuning of text-to-image diffusion models. Our investigation is split into two core aspects: (1) the motivation for adjusting the direction (i.e., angle) of neurons during finetuning, and (2) the reasoning for utilizing orthogonal transformations specifically to carry out this adjustment.

Addressing the first question, we are motivated by empirical findings reported in \cite{liu2018decoupled,chen2020angular}, which indicate that angular distinctions in features effectively reflect the underlying semantic gap. SphereNet~\cite{liu2017deep} demonstrates that constraining all neural activations to lie on a unit hypersphere results in similar or improved generalization and representational power, suggesting that neuronal directions encapsulate critical data information. To further highlight the role of angular representation, we present a controlled experiment as illustrated in Figure~\ref{fig:toy_ae}. Here, a basic convolutional autoencoder is trained on a set of flower images using the standard inner product for feature extraction ($z$ representing the convolution output for filter $\bm{w}$ and input patch $\bm{x}$). During evaluation, we contrast three feature computation methods: (a) the conventional inner product (consistent with training), (b) magnitude-based features, and (c) angular-based features. Analysis of Figure~\ref{fig:toy_ae} indicates that recovering images using angular information of neurons leads to nearly perfect reconstruction, while neuron magnitudes provide little to no relevant content. Notably, only the inner product is utilized in training, with no exposure to cosine activations in (c). These observations support the notion that semantic information is primarily encoded in neuron directions, and that manipulating these directions is potentially the most effective means of altering an image’s semantics.

In response to the second question, a straightforward technique for adjusting neuron directions is to apply a uniform rotation or reflection to all neurons within a given layer, corresponding naturally to an orthogonal transformation. Although one could adopt alternative angular manipulations allowing each neuron a distinct rotation, we determine that orthogonal transformations offer an optimal balance between adaptability and structural regularity. Further, it is established in \cite{liu2021orthogonal} that orthogonal transformations afford sufficient expressive capacity for neural network training. To corroborate these claims, we conduct a regularization study employing the ControlNet~\cite{zhang2023adding} architecture as detailed in Figure~\ref{fig:ortho}. The outcomes reveal that our canonical OFT, equipped with the orthogonality constraint, achieves reliable and precise control post-training (at epoch 20), while a variant lacking this constraint is unable to synthesize realistic images or demonstrate controllability. These results underscore the critical nature of the orthogonality constraint in the context of OFT.

\subsection{General Framework}

In standard finetuning approaches, models—or specific layers—are updated using gradient descent with a relatively low learning rate. This modest learning rate effectively biases the optimization towards remaining close to the original pretrained parameters. Essentially, conventional finetuning aims to minimize the deviation from the original weights, formalized by $\thickmuskip=2mu \medmuskip=2mu \| \bm{M} - \bm{M}^0\|$, where $\bm{M}$ and $\bm{M}^0$ represent the finetuned and pretrained weights, respectively. While this implicit regularization is intuitively reasonable, it may not sufficiently restrict changes for very large models. To address this, the LoRA technique incorporates a further constraint by enforcing low-rank structure on the weight update: specifically, $\thickmuskip=2mu \medmuskip=2mu \text{rank}(\bm{M} - \bm{M}^0)=r'$, with $r'$ being a small integer. Unlike LoRA, the OFT framework imposes a different form of constraint that targets neuron pairwise similarity via hyperspherical energy: $\thickmuskip=2mu \medmuskip=2mu \| \text{HE}(\bm{M})-\text{HE}(\bm{M}^0) \|=0$, where $\text{HE}(\cdot)$ refers to the hyperspherical energy computed for a weight matrix.

For example, take a fully connected layer defined as $\thickmuskip=2mu \medmuskip=2mu \bm{W}=\{\bm{w}_1,\cdots,\bm{w}_n\}\in\mathbb{R}^{d\times n}$, where each column $\bm{w}_i\in\mathbb{R}^{d}$ is the $i$-th neuron (and $\bm{W}^0$ the pretrained counterpart). The output of this layer, $\thickmuskip=2mu \medmuskip=2mu \bm{z}\in\mathbb{R}^n$, is computed by $\thickmuskip=2mu \medmuskip=2mu \bm{z}=\bm{W}^\top\bm{x}$, with input $\thickmuskip=2mu \medmuskip=2mu \bm{x}\in\mathbb{R}^d$. The OFT paradigm seeks to minimize the difference in hyperspherical energy between the pretrained and adapted models:

\begin{equation}\label{eq:oft_principle}
\footnotesize
\min_{\bm{W}} \left\| \text{HE}(\bm{W})-\text{HE}(\bm{W}^0)\right\|~~\Leftrightarrow~~\min_{\bm{W}} \bigg{\|} \sum_{i\neq j} \|\hat{\bm{w}}_i-\hat{\bm{w}}_j\|^{-1}-\sum_{i\neq j} \|\hat{\bm{w}}^0_i-\hat{\bm{w}}^0_j\|^{-1}\bigg{\|}
\end{equation}

Here, $\thickmuskip=2mu \medmuskip=2mu \hat{\bm{w}}_i:=\bm{w}_i/\|\bm{w}_i\|$ designates the normalized version of neuron $i$, while the hyperspherical energy of $\bm{W}$ is defined as $\thickmuskip=2mu \medmuskip=2mu \text{HE}(\bm{W}):= \sum_{i\neq j} \|\hat{\bm{w}}_i-\hat{\bm{w}}_j\|^{-1}$. Notably, Eq.~\eqref{eq:oft_principle} reaches its minimal value of zero when the finetuned weights $\bm{W}$ differ from the pretrained $\bm{W}^0$ only by an orthogonal transformation—that is, $\thickmuskip=2mu \medmuskip=2mu \bm{W}=\bm{R}\bm{W}^0$, with $\thickmuskip=2mu \medmuskip=2mu \bm{R}\in\mathbb{R}^{d\times d}$ being an orthogonal matrix (determinant $1$ or $-1$, corresponding to rotation or reflection, respectively). OFT leverages this property, optimizing neural networks by learning orthogonal, layer-wide matrices to transform neurons. Specifically, OF

\begin{equation}\label{eq:oft_general}
    \footnotesize
    \bm{z}=\bm{W}^\top\bm{x}=(\bm{R}\cdot\bm{W}^0)^\top\bm{x},~~~\text{s.t.}~\bm{R}^\top\bm{R}=\bm{R}\bm{R}^\top=\bm{I}
\end{equation}

Here, $\bm{W}$ indicates the weight matrix obtained from OFT finetuning, while $\bm{I}$ refers to the identity matrix. An overview of OFT can be found in Figure~\ref{fig:block}. Mirroring the zero initialization scheme found in LoRA, OFT requires initializing the finetuning process such that the pretrained model starts from $\bm{W}^0$. To do so, the orthogonal matrix $\bm{R}$ is set to the identity matrix at initialization, ensuring the finetuned parameters initially coincide with the pretrained weights. Maintaining the orthogonality of $\bm{R}$ throughout training can be accomplished by leveraging differentiable orthogonalization techniques as outlined in \cite{liu2021orthogonal,lezcano2019cheap}. Further discussion on practical and efficient approaches to enforce orthogonality is provided below.

\subsection{Efficient Orthogonal Parameterization}\label{sect:eff_ortho}

Although the Gram-Schmidt process yields differentiable orthogonal matrices, its computational overhead is prohibitive for practical purposes~\cite{liu2021orthogonal}. To mitigate this issue, we employ the Cayley parameterization for constructing orthogonal matrices. Specifically, this entails defining $\bm{R}=(\bm{I}+\bm{Q})(I-\bm{Q})^{-1}$, where $\bm{Q}$ is a skew-symmetric matrix that fulfills $\bm{Q}=-\bm{Q}^\top$. The benefit of this method lies in its computational efficiency, although it restricts generated orthogonal matrices to have determinant $1$, thus limiting $\bm{R}$ to the special orthogonal group. Empirically, we observe that this constraint does not hinder performance. Even so, Cayley-based parameterization can become parameter-intensive for high-dimensional settings (large $d$). To improve efficiency, we suggest decomposing $\bm{R}$ into a block-diagonal structure composed of $r$ blocks, yielding:

\begin{equation}
    \footnotesize
    \bm{R}=\text{diag}(\bm{R}_1,\bm{R}_2,\cdots,\bm{R}_r)=\begin{bmatrix}
    \bm{R}_1\in \text{O}(\frac{d}{r}) &  &\\
    &\ddots & \\
    & & \bm{R}_r\in \text{O}(\frac{d}{r})
    \end{bmatrix}\in \text{O}(d)
\end{equation}

where $\text{O}(d)$ refers to the group of $d$-dimensional orthogonal matrices, with $\thickmuskip=2mu \medmuskip=2mu \bm{R} \in \mathbb{R}^{d \times d}$ and, for each $i$, $\thickmuskip=2mu \medmuskip=2mu \bm{R}_i \in \mathbb{R}^{d/r \times d/r}$ as orthogonal blocks. In the special case where $\thickmuskip=2mu \medmuskip=2mu r = 1$, the resulting block-diagonal orthogonal matrix simply reduces to a standard full orthogonal matrix without additional restrictions. For a general orthogonal matrix of dimension $\thickmuskip=2mu \medmuskip=2mu d \times d$, the total number of free parameters is $\thickmuskip=2mu \medmuskip=2mu d(d-1)/2$, leading to $\mathcal{O}(d^2)$ parameter complexity. When using an $r$-block diagonal structure, the parameter count becomes $\thickmuskip=2mu \medmuskip=2mu d(d/r-1)/2$, corresponding to $\thickmuskip=2mu \medmuskip=2mu \mathcal{O}(d^2/r)$. To further decrease the number of parameters, one can optionally employ block sharing so that $\thickmuskip=2mu \medmuskip=2mu \bm{R}_i = \bm{R}_j$ for all $i \neq j$, which brings the complexity down to $\mathcal{O}(d^2/r^2)$. Despite these various parameter reduction techniques, the composite matrix $\bm{R}$ continues to satisfy the orthogonality condition, ensuring the hyperspherical energy is fully retained.

We now compare the parameter efficiency of OFT with that of LoRA. In LoRA, assuming a low-rank dimension of $r'$, the trainable parameter count amounts to $\thickmuskip=2mu \medmuskip=2mu r'(d+n)$. Considering scenarios where both $r$ and $r'$ scale proportionally with $d$, such as $\thickmuskip=2mu \medmuskip=2mu r = r' = \alpha d$ with some $0 < \alpha \leq 1$, the LoRA method attains a parameter complexity of $\mathcal{O}(d^2 + dn)$, whereas OFT exhibits a much more favorable complexity of $\mathcal{O}(d)$. For instance, suppose the weight matrix has size $\thickmuskip=2mu \medmuskip=2mu 128 \times 128$; with $\thickmuskip=2mu \medmuskip=2mu r' = 8$, LoRA would require $2,048$ learnable parameters, while OFT, using $\thickmuskip=2mu \medmuskip=2mu r=8$ (and without block sharing), would involve $960$ trainable parameters.

\subsection{Constrained Orthogonal Finetuning}

The expressivity of basic OFT can be further reduced by constraining the finetuned parameters to remain close to their pretrained values. In particular, COFT implements the following computation during inference:

\begin{equation}\label{eq:coft_general}
    \footnotesize
    \bm{z} = \bm{W}^\top \bm{x} = (\bm{R} \cdot \bm{W}^0)^\top \bm{x}, ~~~ \text{s.t.}~\bm{R}^\top\bm{R} = \bm{R}\bm{R}^\top = \bm{I}, ~ \norm{\bm{R}-\bm{I}}\leq \epsilon
\end{equation}

This imposes not only an orthogonality constraint but also requires that $\bm{R}$ does not deviate from the identity by more than $\epsilon$. While orthogonality can be enforced using the Cayley transform (see Section~\ref{sect:eff_ortho}), enforcing the $\epsilon$-neighborhood constraint is not straightforward with Cayley-parametrized orthogonal matrices. To better understand the behavior of the Cayley transform, one can apply a Neumann series to approximate $\thickmuskip=2mu \medmuskip=2mu \bm{R} = (\bm{I} + \bm{Q})(\bm{I} -

series converges under the operator norm. Consequently, we can incorporate the constraint $\thickmuskip=2mu \medmuskip=2mu\|\bm{R}-\bm{I}\|\leq\epsilon$ within the Cayley transform, resulting in an equivalent requirement: $\thickmuskip=2mu \medmuskip=2mu \|\bm{Q}-\bm{0}\|\leq\epsilon'$, where $\bm{0}$ denotes the zero matrix and $\epsilon'$ is an alternative tolerance hyperparameter, distinct from $\epsilon$. This constraint on $\bm{Q}$ lends itself readily to enforcement via projected gradient descent. For identity initialization of the orthogonal matrix $\bm{R}$, one simply sets $\bm{Q}$ to the zero matrix. COFT effectively integrates two explicit constraints: limiting the difference in hyperspherical energy and bounding the departure from the pretrained parameters. Although the latter is often only implicit in existing finetuning strategies, COFT formalizes it as an explicit criterion. Empirically, while OFT already exhibits strong performance, the explicit regularization in COFT leads to even greater stability during finetuning. Figure~\ref{fig:eps_coft} illustrates how varying $\epsilon$ influences COFT’s behavior; increasing $\epsilon$ allows the COFT-finetuned model to approach the OFT-finetuned model in behavior, whereas decreasing $\epsilon$ drives the COFT-finetuned model to behave more like the original pretrained text-to-image diffusion model.

\subsection{Re-scaled Orthogonal Finetuning}

We introduce a straightforward modification of OFT in which each neuron’s activation is further modulated by a learned scaling factor. This modification is motivated by the observation that scaling does not alter the hyperspherical energy, as the normalization process removes magnitude effects. More precisely, the forward propagation follows: $\thickmuskip=2mu \medmuskip=2mu\bm{z}=(\bm{R}\bm{W}^0\bm{D})^\top\bm{x}$\footnote{\textbf{Errata}: The camera-ready NeurIPS paper incorrectly stated the re-scaled OFT forward pass as $\thickmuskip=2mu \medmuskip=2mu\bm{z}=(\bm{D}\bm{R}\bm{W}^0)^\top\bm{x}$. However, the code implementation was accurate, so findings in Appendix~\ref{app:rescaled} remain valid.}, where $\thickmuskip=2mu \medmuskip=2mu \bm{D}=\text{diag}(s_1,\cdots,s_n)\in\mathbb{R}^{n\times n}$ is a trainable diagonal matrix, and all $s_1,\cdots,s_n$ are strictly positive. Unlike the original OFT formulation in Eq.~\eqref{eq:oft_general}, where only $\bm{R}$ is adapted, this variant learns both the diagonal matrix $\bm{D}$ and the orthogonal matrix $\bm{R}$. Adding this neuron-wise re-scaling enhances OFT's modeling capability with a negligible increase in parameter count. Although we primarily use standard OFT in our main experiments to isolate the effect of orthogonal transformation, our findings indicate that the re-scaled variant typically achieves improved results (see Appendix~\ref{app:rescaled}).

\section{Intriguing Insights and Discussions}

\textbf{OFT exhibits architecture independence}. The OFT methodology can, in theory, be deployed across various neural network architectures. In the context of Transformers, LoRA modifications generally target the attention weights~\cite{hulora2022}. To ensure fair comparisons with LoRA, our implementation restricts OFT finetuning exclusively to the attention parameters. OFT is also naturally compatible with convolutional layers due to the interpretation of $\bm{R}$'s block-diagonal composition, which offers additional structural meaning in convolutional contexts (unlike LoRA). If the number of blocks equals the count of input channels, each block modifies only its corresponding neuron channel, which closely resembles depth-wise convolutional kernel learning~\cite{chollet2017xception}. Alternatively, when the same block is reused across all channels, an identical orthogonal matrix governs the transformation in each channel. A preliminary assessment of OFT applied to convolutional layers is included in Appendix~\ref{app:convolution}.

\textbf{Connection to LoRA}. LoRA incorporates a low-rank matrix to limit substantial alterations to the information encoded in the pretrained weights. In comparison, OFT ensures the transformation applied to the pretrained weights remains orthogonal (thus full-rank), which also safeguards the integrity of pretraining. The OFT forward computation can be expressed as $\thickmuskip=2mu \medmuskip=2mu \bm{z}=(\bm{R}\bm{W}^0)^\top\bm{x}=(\bm{W}^0+(\bm{R}-\bm{I})\bm{W}^0)^\top\bm{x}$, where the term $\thickmuskip=2mu \medmuskip=2mu (\bm{R}-\bm{I})\bm{W}^0$ plays a similar role to the low-rank adjustment in LoRA. As $\bm{W}^0$ is generally of full rank, OFT will also perform a low-rank update provided $\thickmuskip=2mu \medmuskip=2mu \bm{R}-\bm{I}$ is low-rank. Analogous to the rank parameter $r'$ in LoRA, OFT introduces a diagonal block size $r$ as a means of controlling the count of learnable parameters. Fundamentally, LoRA and OFT showcase two alternative strategies for parameter-efficient adaptation: LoRA leverages the low-rank property to cut down parameter count, while OFT achieves parameter efficiency by leveraging block-diagonal (sparse) orthogonal structure.

\textbf{Why OFT converges faster}. There are two perspectives that elucidate the rapid convergence of OFT. Firstly, Figure~\ref{fig:toy_ae} reveals that the most influential adjustments affecting semantic change stem from rotating neuron directions—a capability intrinsic to OFT. Secondly, the optimization process of OFT can be interpreted as neuron adjustment on a smooth hypersphere manifold, promoting a more favorable optimization landscape. This effect is also supported by empirical observations in \cite{liu2021orthogonal}.

\textbf{Why not minimize hyperspherical energy}. Unlike \cite{liu2021orthogonal}, our objective does not involve minimizing hyperspherical energy. In classification scenarios, redundancy-free neurons are typically desirable, but minimizing hyperspherical energy leads to neuron configurations uniformly distributed over the hypersphere. This arrangement is not suitable for finetuning, as such uniformity may undermine the knowledge acquired during pretraining.

\textbf{Balancing flexibility and regularization in finetuning}. Our findings highlight an intrinsic balance between flexibility and regularity during finetuning. Traditional finetuning techniques offer the greatest adaptability, yet they often result in instability and can readily lead to model collapse. In contrast, OFT, through the simple mechanism of maintaining pairwise neuron angles, strikes a favorable compromise between flexibility and structure. Moreover, the block-diagonal parameterization functions as a more stringent regularizer for the orthogonal transformation.

\textbf{No inference-time computational cost}. Differing from approaches like ControlNet, our OFT method does not impose any additional inference-time computation for the adapted model. During inference, it suffices to directly left-multiply the learned orthogonal matrix $\bm{R}$ with the pretrained weight matrix $\bm{W}^0$, yielding a new effective weight matrix $\thickmuskip=2mu \medmuskip=2mu \bm{W}=\bm{R}\bm{W}^0$. Consequently, inference speed remains on par with the original pretrained model.

\section{Experiments and Results}

\textbf{General experimental setup}. Our experiments utilize Stable Diffusion v1.5~\cite{rombach2022high} as the base pretrained text-to-image generator. For fairness, generated outputs from each method are randomly selected. Subject-driven tasks largely follow the protocol outlined in DreamBooth~\cite{ruiz2023dreambooth}, while methods for controllable generation are aligned with the frameworks of ControlNet~\cite{zhang2023adding} and T2I-Adapter~\cite{mou2023t2i}. To guarantee a proper comparison with LoRA, both OFT and COFT are restricted to application within the same layer as LoRA. Additional implementation specifics and further results appear in Appendix~\ref{app:settings}.

\subsection{Subject-driven Generation}

\textbf{Setup}. DreamBooth~\cite{ruiz2023dreambooth} and LoRA~\cite{hulora2022} are adopted as baseline approaches. Consistent loss functions are employed across all methods as prescribed by DreamBooth. For both DreamBooth and LoRA, we adhere closely to the experimental protocols and recommended hyperparameter selections described in their respective original works. Additional quantitative and qualitative results are presented in Appendices~\ref{app:settings},~\ref{app:coft_oft},~\ref{app:more_qualitative}, and~\ref{app:failure}.

\textbf{Finetuning stability and convergence}. To begin, we assess the stability during finetuning as well as the convergence rates for DreamBooth, LoRA, OFT, and COFT. The corresponding findings are depicted in Figure~\ref{fig:teaser} and Figure~\ref{fig:db_convergence}. It is evident that both COFT and OFT provide robust and stable finetuning for the diffusion models. Within 400 iterations, DreamBooth and the OFT variants are capable of exerting effective control over the generative process, but LoRA fails to adequately maintain the subject’s identity. As training proceeds to 2000 iterations, DreamBooth starts producing degraded, collapsed outputs and LoRA is unable to depict a yellow shirt, instead generating yellow fur. In stark contrast, OFT and COFT continue to exercise reliable and steady control over the generation. These experiments confirm that the OFT framework achieves both rapid and stable convergence in the context of subject-driven generation tasks. It is noteworthy that this robustness to the number of finetuning iterations is crucial, as it eliminates the need for careful adjustment of iteration counts for distinct subjects. With both OFT and COFT, a sufficiently large iteration count may be selected in advance, eliminating the need for fine-tuning. In the case of COFT, with an appropriately chosen $\epsilon$, both the learning rate and iteration number can be specified with minimal effort.

\textbf{Quantitative comparison}. In accordance with \cite{ruiz2023dreambooth}, we perform a quantitative evaluation covering subject fidelity (using DINO~\cite{caron2021emerging} and CLIP-I~\cite{radford2021learning}), text prompt fidelity (via CLIP-T~\cite{radford2021learning}), and image diversity (LPIPS~\cite{zhang2018unreasonable}). The CLIP-I score is computed as the mean pairwise cosine similarity between CLIP embeddings of synthesized and original images. DINO follows the same procedure as CLIP-I but utilizes ViT S/16 DINO representations. CLIP-T measures the average cosine similarity between CLIP embeddings of the text prompt and generated outputs. For diversity assessment, we report the average LPIPS cosine similarity across generated samples of the same subject and prompt. According to Table~\ref{table:dreambooth_final}, COFT and OFT outperform DreamBooth and LoRA by a significant margin in the DINO and CLIP-I metrics, and display marginally better or similar results in prompt fidelity and diversity. For each technique, we repeat finetuning 30 times on the same pretrained model using distinct random seeds, thereby mitigating stochastic effects. The outcomes indicate that the OFT framework not only demonstrates superior convergence and stability but also consistently produces improved final results.

\textbf{Qualitative comparison}. To better grasp the advantages offered by OFT, we present a set of randomly selected samples from the subject-driven generation task in Figure~\ref{fig:dreambooth_final}. For consistency, all samples are produced using the identical finetuned model across different methods, ensuring that text prompts are not individually optimized for each output. For every approach, the version of the model that records the highest validation CLIP score is chosen. Examining Figure~\ref{fig:dreambooth_final}, it is apparent that OFT and COFT are both highly effective in maintaining the semantic features of the subject, whereas LoRA frequently does not retain the subject's identity (\emph{e.g.}, LoRA entirely misses the subject in the bowl case). Furthermore, both OFT and COFT display superior accuracy in leveraging text prompts to guide image generation; by contrast, while DreamBooth is capable of keeping the subject identity, it often fails to render the image as described by the textual prompt (\emph{e.g.}, the initial row of the bowl illustration). This comparative qualitative analysis highlights that the OFT framework provides improved controllability over the image generation process alongside robust subject preservation. Additionally, OFT demonstrates robustness to the iteration count—showing stable performance across a wide range of iterations—unlike DreamBooth and LoRA, which are more sensitive to this factor. Further qualitative illustrations are available in Appendix~\ref{app:more_qualitative}. We also include a human study in Appendix~\ref{app:human}, which substantiates the advantages of OFT.

\subsection{Controllable Generation}

\textbf{Settings}. As baselines, we employ ControlNet~\cite{zhang2023adding}, T2I-Adapter~\cite{mou2023t2i}, and LoRA~\cite{hulora2022}. In the main manuscript, we focus on three challenging tasks requiring controllable generation: Canny edge to image~(C2I) using the COCO dataset~\cite{lin2014microsoft}, segmentation map to image~(S2I) on the ADE20K dataset~\cite{zhou2017scene}, and landmark to face~(L2F) leveraging the CelebA-HQ dataset~\cite{karras2018progressive,xia2021tedigan}. All techniques are employed to finetune Stable Diffusion~(SD) v1.5 for 20 training epochs on each dataset. Additional findings are provided in Appendix~\ref{app:more_qualitative}, \ref{app:more_control_task}, \ref{app:failure}.

\textbf{Convergence}. We assess how quickly ControlNet, T2I-Adapter, LoRA, and COFT converge on the L2F task, using both numerical metrics and visual comparisons. For our quantitative measure, we calculate the mean $\ell_2$ distance between the control face landmarks and the model-predicted landmarks. Figure~\ref{fig:face_convergence} illustrates how face landmark errors evolve as models are finetuned over varying numbers of epochs. The results indicate that COFT and OFT reach optimal performance in far fewer epochs compared to the other approaches. For instance, LoRA requires 20 epochs to match the performance that OFT attains at just the 8-th epoch. Importantly, both OFT and COFT involve approximately the same number of trainable parameters as LoRA—which is considerably less than ControlNet—yet their convergence is markedly faster than previous methods. Additional support for OFT's rapid convergence is found in Figure~\ref{fig:teaser}. Specifically, the example on the right in Figure~\ref{fig:teaser} demonstrates OFT's superior data efficiency: it achieves effective convergence with merely 5\% of the ADE20K dataset, outperforming both ControlNet and LoRA in this respect. In our qualitative analysis, we center our comparisons on OFT, COFT, and ControlNet, as ControlNet delivers the next-best landmark accuracy to OFT. According to Figure~\ref{fig:face_convergence_qual}, both OFT and COFT exhibit steady convergence; the synthesized face poses incrementally align with the reference landmarks. In contrast, ControlNet displays an abrupt improvement in controllability, which only materializes after the 8-th epoch, corroborating the sudden convergence described in \cite{zhang2023adding}. The smooth and stable convergence of OFT is especially advantageous for practical use, as it yields more consistent training behavior, facilitating the search for suitable OFT hyperparameters.

\textbf{Quantitative comparison}. To assess controllable generation performance, we propose a control consistency metric. This metric works by extracting the control signal from each generated image and comparing it to the original input control. Specifically, for the C2I task, we measure IoU and F1 scores, while for S2I, we report mean IoU, as well as mean and overall accuracy. For L2F, we calculate the average $\ell_2$ distance between the input control landmarks and those predicted by the model. Additional technical details on these metrics are described in Appendix~\ref{app:settings}. Each baseline and proposed method is evaluated using its optimal hyperparameters. As shown in Table~\ref{table:control_gen}, both OFT and COFT significantly outperform other baselines in delivering strong and precise control. We find that adapter-based techniques (such as T2I-Adapter and ControlNet) show slower convergence rates and tend to produce inferior outcomes. Although LoRA exceeds ControlNet in S2I performance, it falls short in C2I and L2F. Overall, LoRA's maximum performance appears limited, even when its hyperparameters are exhaustively tuned. In contrast, our experiments indicate that the OFT framework continues to improve as the number of trainable parameters increases, with no clear performance saturation observed. Notably, our approach to quantitatively evaluating control in generation represents a key innovation, as it enables rigorous assessment of control fidelity for finetuned models, bounded only by the segmentation/detection model’s own accuracy.

\textbf{Qualitative comparison}. We conduct a qualitative evaluation of OFT and COFT in relation to leading approaches such as ControlNet, T2I-Adapter, and LoRA. As illustrated by the randomly sampled images in Figure~\ref{fig:control_final}, OFT and COFT are capable of generating images that are both realistic and of high fidelity, while simultaneously offering precise controllability. For the S2I task, it is evident that LoRA fails to follow the guidance of the input segmentation map, whereas ControlNet, OFT, and COFT successfully maintain effective control over image synthesis. Compared to ControlNet, OFT and COFT produce visually superior images, exhibiting richer detail and more plausible geometric consistency, while utilizing substantially fewer model parameters. Regarding the C2I task, OFT and COFT can generate semantically coherent images based on coarse Canny edge inputs, outperforming T2I-Adapter and LoRA, which achieve inferior results. In the L2F task, our framework demonstrates the highest accuracy in controlling facial poses, even when presented with challenging inputs. Across all three control scenarios, OFT and COFT consistently deliver higher-quality images than existing state-of-the-art models, thereby validating the efficacy of our OFT approach for controllable image generation. For a broader perspective on qualitative performance, we provide additional visual results for the three tasks in Appendix~\ref{app:control_exp}, and further exhibit OFT’s versatility on supplementary control tasks—including dense pose-to-body, sketch-to-image, and depth-to-image translation—in Appendix~\ref{app:more_control_task}.

\section{Concluding Remarks and Open Problems}

Building on the insight that angular relationships among neurons are key to shaping visual semantics, this work introduces a straightforward yet powerful finetuning approach—orthogonal finetuning (OFT)—designed to control text-to-image diffusion models. The approach is applied in two contexts: subject-driven generation and controllable generation. In comparison with current techniques, OFT achieves enhanced controllability and greater stability during finetuning, all while employing a reduced set of parameters. Furthermore, OFT does not add inference-time computational costs, enabling practical and efficient deployment.

The introduction of OFT also opens several notable avenues for future investigation. Firstly, the orthogonality constraint in OFT is ensured using Cayley parametrization, which necessitates computing a matrix inverse—a factor that slightly constrains the scalability of the method. While this is partially mitigated via block diagonal parametrization, developing a differentiable and more efficient computation for this matrix inverse remains unresolved. Secondly, OFT presents a unique advantage in terms of compositionality: the orthogonal matrices yielded by independent OFT finetuning tasks can be multiplied together, with their product still being orthogonal. Whether this multiplicative structure retains information across all corresponding downstream tasks warrants further exploration. Lastly, the method's parameter efficiency is tied closely to the block diagonal construction, which may inadvertently introduce biases and curtail flexibility. Pursuing strategies to enhance parameter efficiency in a less biased and more effective fashion stands out as a pressing challenge for future research.

\end{document}